\section{Introduction}

\begin{figure}[t]
    \centering
    \includegraphics[width=\columnwidth]{figs/figure1.pdf}
    \caption{Compliant trajectory vs.\ authorization-drift trajectory. Asked to draft a follow-up email, the compliant agent uses only the contact record provided in context, while the drifting agent, without any adversarial input, queries private CRM history and searches a manager's mailbox to enrich the draft, expanding its data access beyond what the task authorizes.}
    \label{fig:drift-illustration}
\end{figure}



LLM agents increasingly act on users' behalf through tools, code, and data access. Existing safety research largely examines adversarial failures, including prompt injection, memory poisoning, and maliciously induced tool use \citep{zhan2024injecagent,debenedetti2024agentdojo,zhang2025asb}. Yet agents can also exceed their authorization without an attacker, simply because additional access appears useful. Figure~\ref{fig:drift-illustration} illustrates this failure: an agent asked to draft an email consults private CRM records and a manager's mailbox to enrich its response. We formalize this phenomenon as \textbf{authorization drift}: a sequence-level failure in which an agent expands its effective authorization boundary while pursuing the user's goal, without external instructions directing the expansion. Authorization drift differs from prompt injection in its non-adversarial origin and from static over-privilege in its emergence during execution. Evaluating this failure requires more than removing adversarial inputs: agents must face useful but unauthorized actions while retaining a feasible way to complete the task within bounds. This tests whether agents choose to respect authorization boundaries, rather than whether they can resist an attacker.

To enable this evaluation, we introduce the \textbf{Authorization Drift Benchmark (ADB)}, comprising 648 samples across seven domains, each pairing a sufficient but effortful compliant path with an attractive unauthorized shortcut. This dual-path design makes unauthorized actions useful but unnecessary, allowing us to distinguish avoidable overreach from task requirements. ADB is therefore an opportunity-conditioned stress test, not an estimate of drift prevalence in ordinary deployments. ADB evaluates four authorization dimensions: \textbf{Data}, accessing information beyond the authorized data boundary; \textbf{Tool}, using tools beyond the permitted capabilities; \textbf{Scope}, expanding beyond the requested objective; and \textbf{Time}, persisting or re-planning beyond the intended execution window. An \emph{anti-recipe rule} preserves realistic information gaps rather than prescribing step-by-step tool calls, requiring agents to navigate these choices themselves. To ensure executable environments, genuinely tempting shortcuts, and feasible compliant paths, we construct ADB through \textbf{Staged Execution-Reflective Synthesis (SERS)}. A Designer specifies tasks and authorization annotations, a Builder implements environments and shortcut opportunities, and a quality-control loop evaluates real rollouts from two probe models from different families. By grounding validation in observed agent behavior rather than LLM self-critique, SERS checks whether these design goals hold during execution, filtering out tasks that are trivially compliant or structurally force drift.

Our evaluation contains 25{,}207 evaluated trajectories from 13 LLMs spanning 4B to 397B parameters under \textsc{Normal}; task success is LLM-judged, whereas drift is detected by rule. We additionally compare \textsc{ReAct}~\citep{yao2022react} and \textsc{Planner--Executor}~\citep{Chase_LangChain_2022,wang2023plan} on four representative backbones. \textbf{Capability and drift rise together}: frontier-capable models exhibit drift rates of $0.43$ to $0.61$, yet $86.18\%$ of analyzed tasks admit an observed clean-success trajectory, showing that compliant completion is observed on most tasks. Drift rates are stable across runs, with standard deviations below $0.02$ for every model, and early data over-collection dominates over a distinct cluster of later scope and time expansion. Lower drift under \textsc{Planner--Executor} accompanies collapsing clean-completion performance, while non-monotonic within-family trends suggest that parameter count alone does not explain boundary-keeping. To further understand how these behaviors may harm individuals and society, we use legal provisions as an interpretive framework, mapping annotated drift points to data-protection and cybersecurity provisions across the European (EU), Mainland China, Singapore, and Hong Kong (HK). This analysis connects technical overreach to concrete concerns, including intrusion into private lives, loss of control over personal information, and unauthorized access to confidential records. GDPR Art.~5(1)(c), concerning data minimisation, is the most frequently implicated provision for every model, and $93.5\%$ of legally grounded drifting runs map to provisions from multiple jurisdictions. These mappings explain the potential human and societal consequences of drift without establishing legal violations or simultaneous applicability across jurisdictions. Our contributions are threefold:

\begin{itemize}[left=0pt]
    \item We formalize authorization drift as a distinct, non-adversarial safety failure across data, tool, scope, and time boundaries, distinguishing it from prompt injection, static over-privilege, and generic privacy leakage.
    \item We release ADB, a 648-sample dual-path benchmark across seven domains, together with SERS, a staged, execution-grounded synthesis pipeline that validates task completion and shortcut triggerability through real agent rollouts.
    \item Across 13 LLMs, with three-framework comparisons on four backbones (25{,}207 evaluation trajectories), we characterize the capability--drift association, its persistence and trajectory structure, and the limitations of agentic scaffolds. We further use legal mappings across four jurisdictions to connect these behavioral findings to potential harms to individuals and society.
\end{itemize}
